% BEGIN PROFESSOR FEEDBACK PF27: Discussion - Behaviour categories
\subsection{Differences between behaviour categories}

The category analysis uses motif fingerprints alone. Its class-share and
permutation references establish whether the fingerprints contain category
information. They do not establish whether this representation is more useful
than conventional behavioural features for category classification, which would
require an additional baseline.

RQ3 shows that the Real-CATS categories are not identical in their motif
fingerprints. At the shared two-week window, the model reaches a macro AP of
0.1953 compared with an average class share of 0.1111. The permutation test
finds the category signal under every timestamp ordering. This provides evidence
that the motif distributions contain information about the recorded behaviour
category.

The practical separation is much weaker than the permutation result alone might
suggest. Macro F1 is 0.1914, and the confusion matrix shows that many addresses
are assigned to the wrong category. The one-vs-rest AP values also need to be
read together with their class shares. Phishing Scam has the highest AP, but it
also represents more than two thirds of the sample. Several smaller categories
have only 25 to 47 addresses in the test set, which makes their individual
results less precise.

One likely reason for the limited separation is that a category label does not
define one transaction pattern. Two phishing addresses can use different
collection and withdrawal strategies, while an impersonation scam and a fake
returns scam can produce similar transfers despite having different social
engineering methods. The label describes the purpose of the malicious activity,
whereas the motif fingerprint describes local transaction structure. These two
descriptions are related, but they are not equivalent. The slightly higher AP of
the count-based model suggests that motif volume may also contain category
information. However, the comparison changes more than volume alone and does not
identify which part of the count representation produces the difference.

RQ4 further limits the interpretation of the category results. The development
data selects six different category-specific windows, ranging from six hours to
30 days. On the held-out test set, using these windows instead of the shared
two-week window changes macro AP by only 0.0009. The confidence interval includes
zero and none of the category-level differences is significant after correction
for multiple comparisons. Most selected windows also change when transactions
with the same timestamp are ordered differently.

The development maxima should therefore not be treated as stable properties of
the categories. Selecting the highest point from nine windows will produce a
maximum for every category even when several windows perform almost equally.
The held-out comparison tests whether these apparent differences generalise, and
in this study they do not. The category-specific windows also do not estimate how
long the corresponding scam types last. They only describe where the classifier
performed best on one development sample.
% END PROFESSOR FEEDBACK PF27

